% Options for packages loaded elsewhere
\PassOptionsToPackage{unicode}{hyperref}
\PassOptionsToPackage{hyphens}{url}
\documentclass[
]{article}
\usepackage{xcolor}
\usepackage[margin=1in]{geometry}
\usepackage{amsmath,amssymb}
\setcounter{secnumdepth}{5}
\usepackage{iftex}
\ifPDFTeX
  \usepackage[T1]{fontenc}
  \usepackage[utf8]{inputenc}
  \usepackage{textcomp} % provide euro and other symbols
\else % if luatex or xetex
  \usepackage{unicode-math} % this also loads fontspec
  \defaultfontfeatures{Scale=MatchLowercase}
  \defaultfontfeatures[\rmfamily]{Ligatures=TeX,Scale=1}
\fi
\usepackage{lmodern}
\ifPDFTeX\else
  % xetex/luatex font selection
\fi
% Use upquote if available, for straight quotes in verbatim environments
\IfFileExists{upquote.sty}{\usepackage{upquote}}{}
\IfFileExists{microtype.sty}{% use microtype if available
  \usepackage[]{microtype}
  \UseMicrotypeSet[protrusion]{basicmath} % disable protrusion for tt fonts
}{}
\makeatletter
\@ifundefined{KOMAClassName}{% if non-KOMA class
  \IfFileExists{parskip.sty}{%
    \usepackage{parskip}
  }{% else
    \setlength{\parindent}{0pt}
    \setlength{\parskip}{6pt plus 2pt minus 1pt}}
}{% if KOMA class
  \KOMAoptions{parskip=half}}
\makeatother
\usepackage{graphicx}
\makeatletter
\newsavebox\pandoc@box
\newcommand*\pandocbounded[1]{% scales image to fit in text height/width
  \sbox\pandoc@box{#1}%
  \Gscale@div\@tempa{\textheight}{\dimexpr\ht\pandoc@box+\dp\pandoc@box\relax}%
  \Gscale@div\@tempb{\linewidth}{\wd\pandoc@box}%
  \ifdim\@tempb\p@<\@tempa\p@\let\@tempa\@tempb\fi% select the smaller of both
  \ifdim\@tempa\p@<\p@\scalebox{\@tempa}{\usebox\pandoc@box}%
  \else\usebox{\pandoc@box}%
  \fi%
}
% Set default figure placement to htbp
\def\fps@figure{htbp}
\makeatother
\setlength{\emergencystretch}{3em} % prevent overfull lines
\providecommand{\tightlist}{%
  \setlength{\itemsep}{0pt}\setlength{\parskip}{0pt}}
\usepackage{float}
\usepackage{lscape}
\usepackage{booktabs}
\usepackage{longtable}
\usepackage{array}
\usepackage{multirow}
\usepackage{wrapfig}
\usepackage{colortbl}
\usepackage{pdflscape}
\usepackage{tabu}
\usepackage{threeparttable}
\usepackage{threeparttablex}
\usepackage[normalem]{ulem}
\usepackage{makecell}
\usepackage{xcolor}
\usepackage{bookmark}
\IfFileExists{xurl.sty}{\usepackage{xurl}}{} % add URL line breaks if available
\urlstyle{same}
\hypersetup{
  pdftitle={Analyse des inégalités éducatives au Nigeria},
  pdfauthor={Licence en Économie},
  hidelinks,
  pdfcreator={LaTeX via pandoc}}

\title{Analyse des inégalités éducatives au Nigeria}
\usepackage{etoolbox}
\makeatletter
\providecommand{\subtitle}[1]{% add subtitle to \maketitle
  \apptocmd{\@title}{\par {\large #1 \par}}{}{}
}
\makeatother
\subtitle{Nigeria General Household Survey -- Panel Wave 4}
\author{Licence en Économie}
\date{17 March 2026}

\begin{document}
\maketitle

{
\setcounter{tocdepth}{2}
\tableofcontents
}
\begin{center}\rule{0.5\linewidth}{0.5pt}\end{center}

\section{Introduction}\label{introduction}

L'éducation constitue l'un des piliers fondamentaux du développement
économique. Dans la théorie du capital humain, initiée par Becker (1964)
et Schultz (1961), l'investissement en éducation améliore la
productivité individuelle, stimule la croissance économique et réduit
les inégalités de revenus. Pour un pays comme le Nigeria --- première
économie d'Afrique subsaharienne et nation la plus peuplée du continent
--- la question de l'accès à l'éducation revêt une importance
stratégique majeure.

Malgré des progrès notables depuis l'indépendance en 1960, le Nigeria
continue de faire face à des inégalités éducatives profondes,
structurées par le genre, l'âge, le milieu de résidence et la
localisation géographique. Ce rapport propose une analyse quantitative
de ces inégalités à partir des données du \textbf{Nigeria General
Household Survey (GHS) -- Panel Wave 4}.

\begin{center}\rule{0.5\linewidth}{0.5pt}\end{center}

\section{Contexte éducatif du
Nigeria}\label{contexte-uxe9ducatif-du-nigeria}

\subsection{Structure du système
éducatif}\label{structure-du-systuxe8me-uxe9ducatif}

Le système éducatif nigérian est organisé selon le schéma
\textbf{6-3-3-4} : six ans d'enseignement primaire, trois ans de collège
(\emph{Junior Secondary School}), trois ans de lycée (\emph{Senior
Secondary School}), et quatre ans d'enseignement supérieur.
L'enseignement primaire est officiellement obligatoire depuis
l'Universal Basic Education Act de 2004, qui étend cette obligation aux
neuf premières années de scolarité.

\subsection{Engagements
internationaux}\label{engagements-internationaux}

Le Nigeria s'est engagé dans le cadre de l'\textbf{ODD 4} visant à
assurer l'accès de tous à une éducation de qualité. Les analyses qui
suivent permettent d'évaluer, à partir de données récentes, la distance
qui sépare encore le pays de cet objectif.

\begin{center}\rule{0.5\linewidth}{0.5pt}\end{center}

\section{Données et méthodologie}\label{donnuxe9es-et-muxe9thodologie}

\subsection{Source des données}\label{source-des-donnuxe9es}

Les données proviennent du \textbf{Nigeria General Household Survey
(GHS) -- Panel Wave 4}, collecté par le National Bureau of Statistics
(NBS) en partenariat avec la Banque mondiale. Deux sections sont
mobilisées : sect1\_harvestw4 et sect2\_harvestw4, portant
respectivement sur la démographie et le niveau d'instruction.

\subsection{Construction de la base
analytique}\label{construction-de-la-base-analytique}

Les deux sections sont fusionnées par un \textbf{full join} sur les
identifiants \texttt{hhid} et \texttt{indiv}, afin de conserver
l'intégralité des observations.

\subsection{Variable d'intérêt : niveau
d'éducation}\label{variable-dintuxe9ruxeat-niveau-duxe9ducation}

La variable \texttt{s2aq10} est recodée en cinq catégories ordonnées
:Aucun, Primaire, Junior Secondary, Senior Secondary, Tertiaire. Note :
D'après l'OCDE, le niveau d'éducation est spécifique au diploe obtenu.
Par conséquent, nous l'avons défini sur la base de la variable s2aq10.
Il est donc possible qu'il y ait des individus qui sont scolarisé mais
n'ayant aucun diplome; des enfants par exemple, ces individus seront
considéré comme ayant aucun niveau d'instructions.

\subsection{Méthodes statistiques}\label{muxe9thodes-statistiques}

Les méthodes suivantes sont mobilisées, avec un seuil de significativité
\textbf{α = 0,05} :

\begin{itemize}
\tightlist
\item
  \textbf{Statistiques descriptives} : fréquences et proportions,
  tableaux de contingence.
\item
  \textbf{Test du chi-deux de Pearson} : indépendance entre deux
  variables catégorielles.
\item
  \textbf{V de Cramer} : intensité de l'association entre deux variables
  catégorielles (0 = aucune, 1 = parfaite).
\item
  \textbf{Test de Kruskal-Wallis} : Le test de Kruskal-Wallis (d'après
  William Kruskal et Wilson Allen Wallis), aussi appelé ANOVA
  unidirectionnelle sur rangs (ou ANOVA à un facteur contrôlé sur rangs)
  est une méthode non paramétrique utilisée pour tester si des
  échantillons trouvent leur origine dans la même distribution. Ce test
  s'intéresse aux médianes de k populations (k⩾3) (ou treatment dans la
  littérature en anglais) et propose comme hypothèse nulle que les k
  échantillons sont confondus et proviennent d'un même échantillon
  (combiné) d'une population.
\item
  \textbf{Test post-hoc de Dunn} (correction de Bonferroni) :
  comparaisons par paires après Kruskal-Wallis significatif.
\item
  \textbf{Intervalles de confiance à 95\%} : approximation normale pour
  les proportions.
\end{itemize}

\begin{center}\rule{0.5\linewidth}{0.5pt}\end{center}

\section{Distribution du niveau
d'éducation}\label{distribution-du-niveau-duxe9ducation}

\subsection{Tableau des fréquences}\label{tableau-des-fruxe9quences}

\begin{longtable}[t]{lrr}
\caption{\label{tab:tableau-freq}Distribution du niveau d'éducation dans l'échantillon}\\
\toprule
Niveau d'éducation & Effectif (n) & Proportion (\%)\\
\midrule
\cellcolor{gray!10}{Aucun} & \cellcolor{gray!10}{7628} & \cellcolor{gray!10}{39.6}\\
Primaire & 3975 & 20.6\\
\cellcolor{gray!10}{Senior Secondary} & \cellcolor{gray!10}{3755} & \cellcolor{gray!10}{19.5}\\
Tertiaire & 1975 & 10.2\\
\cellcolor{gray!10}{Junior Secondary} & \cellcolor{gray!10}{1938} & \cellcolor{gray!10}{10.1}\\
\bottomrule
\end{longtable}

\subsection{Visualisation}\label{visualisation}

\begin{figure}[H]

{\centering \includegraphics{rapport_education_nigeria_v2_files/figure-latex/barplot-educ-1} 

}

\caption{Figure 1 — Répartition de la population par niveau d'éducation}\label{fig:barplot-educ}
\end{figure}

\subsection{Interprétation}\label{interpruxe9tation}

La distribution du niveau d'éducation reflète les caractéristiques
structurelles du système éducatif nigérian. La part d'individus sans
instruction formelle illustre les déficits historiques d'accès à
l'éducation, particulièrement pour les générations plus âgées. La
pyramide éducative --- avec un poids décroissant à mesure que l'on
progresse vers les niveaux supérieurs --- est cohérente avec les
phénomènes d'abandon scolaire et de sélection qui caractérisent les
systèmes éducatifs des pays à revenu intermédiaire.

\begin{center}\rule{0.5\linewidth}{0.5pt}\end{center}

\section{Inégalités éducatives selon le
sexe}\label{inuxe9galituxe9s-uxe9ducatives-selon-le-sexe}

\subsection{Tableau de contingence}\label{tableau-de-contingence}

\begin{longtable}[t]{lrrrr}
\caption{\label{tab:tableau-contingence}Tableau de contingence : effectifs par sexe et niveau d'éducation}\\
\toprule
 & Junior.Secondary & Primaire & Senior.Secondary & Tertiaire\\
\midrule
\cellcolor{gray!10}{Homme} & \cellcolor{gray!10}{494} & \cellcolor{gray!10}{1087} & \cellcolor{gray!10}{1994} & \cellcolor{gray!10}{1146}\\
Femme & 473 & 1202 & 1514 & 803\\
\bottomrule
\end{longtable}

\subsection{Visualisation}\label{visualisation-1}

\begin{figure}[H]

{\centering \includegraphics{rapport_education_nigeria_v2_files/figure-latex/barplot-sexe-1} 

}

\caption{Figure 2 — Distribution du niveau d'éducation par sexe (adultes 18+)}\label{fig:barplot-sexe}
\end{figure}

\subsection{Test du chi-deux et V de
Cramér}\label{test-du-chi-deux-et-v-de-cramuxe9r}

\begin{longtable}[t]{ll}
\caption{\label{tab:resultats-chi2}Résultats du test du chi-deux et V de Cramér (sexe × niveau d'éducation)}\\
\toprule
Statistique & Valeur\\
\midrule
\cellcolor{gray!10}{Chi-deux (χ²)} & \cellcolor{gray!10}{71.784}\\
Degrés de liberté & 3\\
\cellcolor{gray!10}{p-valeur} & \cellcolor{gray!10}{1.77e-15}\\
V de Cramér & 0.091\\
\bottomrule
\end{longtable}

\subsection{Interprétation}\label{interpruxe9tation-1}

Le test du chi-deux est \textbf{significatif} (χ² = 71.78, p =
1.77e-15). Le V de Cramér (0.091) indique une association
\textbf{négligeable} entre le sexe et le niveau d'éducation.

Les inégalités de genre dans l'éducation au Nigeria s'expliquent par une
combinaison de facteurs : normes socioculturelles favorisant la priorité
aux garçons, mariage et maternité précoces pour les filles, et
contraintes économiques conduisant à des arbitrages défavorables aux
filles dans les zones rurales et du Nord.

\begin{center}\rule{0.5\linewidth}{0.5pt}\end{center}

\section{Inégalités éducatives selon
l'âge}\label{inuxe9galituxe9s-uxe9ducatives-selon-luxe2ge}

\subsection{Construction des groupes
d'âge}\label{construction-des-groupes-duxe2ge}

Dans cette partie spécifiquement, nous avons fait supprimé les valeurs
manquantes car elles sont incompatibles avec les tests qui vont suivre.

\subsection{Visualisation}\label{visualisation-2}

\begin{figure}[H]

{\centering \includegraphics{rapport_education_nigeria_v2_files/figure-latex/boxplot-age-1} 

}

\caption{Figure 3 — Distribution du niveau d'éducation par groupe d'âge (adultes 18+)}\label{fig:boxplot-age}
\end{figure}

Les seules valeurs abérantes à considérer sont dans le groupe des
``18-30''. \#\# Test de Kruskal-Wallis

Le test de Kruskal-Wallis est hautement significatif : \textbf{X² =
287.28}, df = 3, \textbf{p \textless2e-16}. Il existe des différences
statistiquement significatives de niveau d'éducation entre au moins deux
groupes d'âge.

\subsection{Test post-hoc de Dunn (correction de
Bonferroni)}\label{test-post-hoc-de-dunn-correction-de-bonferroni}

\begin{longtable}[t]{lrll}
\caption{\label{tab:tableau-dunn}Test post-hoc de Dunn avec correction de Bonferroni}\\
\toprule
Comparaison & Statistique Z & p-valeur ajustée & Significatif\\
\midrule
\cellcolor{gray!10}{18-30 - 31-45} & \cellcolor{gray!10}{4.594} & \cellcolor{gray!10}{1.30e-05} & \cellcolor{gray!10}{Oui *}\\
18-30 - 46-60 & 9.129 & < 2e-16 & Oui *\\
\cellcolor{gray!10}{31-45 - 46-60} & \cellcolor{gray!10}{4.956} & \cellcolor{gray!10}{2.16e-06} & \cellcolor{gray!10}{Oui *}\\
18-30 - 60+ & 15.844 & < 2e-16 & Oui *\\
\cellcolor{gray!10}{31-45 - 60+} & \cellcolor{gray!10}{12.598} & \cellcolor{gray!10}{< 2e-16} & \cellcolor{gray!10}{Oui *}\\
\addlinespace
46-60 - 60+ & 8.295 & 3.26e-16 & Oui *\\
\bottomrule
\end{longtable}

\subsection{Interprétation}\label{interpruxe9tation-2}

Toutes les paires de groupes d'âge diffèrent significativement. La
statistique Z croît avec l'écart entre les groupes : la différence la
plus prononcée oppose les \textbf{18-30 ans} aux \textbf{60 ans et
plus}, confirmant un fort \textbf{effet de génération}. Ce gradient
reflète l'expansion du système éducatif nigérian depuis les années 1970.
Du point de vue économique, ce résultat est cohérent avec la théorie du
capital humain : les investissements publics en éducation ont produit
des effets mesurables sur le niveau d'instruction des cohortes les plus
récentes.

\begin{center}\rule{0.5\linewidth}{0.5pt}\end{center}

\section{Scolarisation des 6--17 ans selon le milieu de
résidence}\label{scolarisation-des-617-ans-selon-le-milieu-de-ruxe9sidence}

\subsection{Tableau de contingence}\label{tableau-de-contingence-1}

\begin{longtable}[t]{lrr}
\caption{\label{tab:tableau-scol}Taux de scolarisation (%) par milieu de résidence (6-17 ans)}\\
\toprule
 & Scolarisé & Non scolarisé\\
\midrule
\cellcolor{gray!10}{Urbain} & \cellcolor{gray!10}{92.3} & \cellcolor{gray!10}{7.7}\\
Rural & 88.5 & 11.5\\
\bottomrule
\end{longtable}

Une majorité de la population entre 6 et 17 ans sont scolarisé dans les
milieux urbains comme ruraux. Cependant, la proportion de jeunes non
scolarisé est plus forte dans le milieu rural. \textbf{Test du chi-deux
:} χ² = 24.63, df = 1, p = 6.94e-07.

\subsection{Visualisation}\label{visualisation-3}

\begin{figure}[H]

{\centering \includegraphics{rapport_education_nigeria_v2_files/figure-latex/barplot-scolar-1} 

}

\caption{Figure 4 — Taux de scolarisation (6-17 ans) par milieu de résidence avec IC à 95\%}\label{fig:barplot-scolar}
\end{figure}

\subsection{Interprétation}\label{interpruxe9tation-3}

L'écart de scolarisation entre milieu urbain et rural s'explique par des
mécanismes agissant à la fois sur l'offre et la demande d'éducation. Du
côté de l'offre, les zones rurales sont moins bien dotées en
infrastructures scolaires. Du côté de la demande, le coût d'opportunité
de la scolarisation est plus élevé en milieu rural, où les enfants
participent davantage aux activités agricoles familiales, et les revenus
des ménages sont en moyenne plus faibles.

\begin{center}\rule{0.5\linewidth}{0.5pt}\end{center}

\section{Disparités géographiques entre
États}\label{disparituxe9s-guxe9ographiques-entre-uxe9tats}

\subsection{Visualisation}\label{visualisation-4}

\newpage
\begin{landscape}

\begin{figure}[H]

{\centering \includegraphics[width=1\linewidth]{rapport_education_nigeria_v2_files/figure-latex/heatmap-1} 

}

\caption{Figure 5 — Part d'adultes sans instruction par État nigérian (quintiles)}\label{fig:heatmap}
\end{figure}

\end{landscape}
\newpage

\subsection{Interprétation}\label{interpruxe9tation-4}

Les disparités géographiques en matière d'éducation reflètent de
profondes inégalités régionales de développement économique. La fracture
Nord-Sud constitue la ligne de clivage principale : les États du
Nord-Ouest et du Nord-Est --- régions moins intégrées aux circuits
économiques formels et affectées par l'insécurité --- tendent à
concentrer les déficits éducatifs les plus importants. À l'inverse, les
États du Sud-Ouest (Lagos, Ogun, Oyo) et du Sud-Est bénéficient d'une
tradition plus ancienne de scolarisation. Ces disparités contribuent à
la persistance des inégalités régionales de productivité et justifient
des politiques ciblées de rattrapage éducatif.

\begin{center}\rule{0.5\linewidth}{0.5pt}\end{center}

\section{Limites méthodologiques}\label{limites-muxe9thodologiques}

\textbf{Causalité vs.~corrélation.} Les associations mises en évidence
sont de nature descriptive. L'identification de mécanismes causaux
nécessiterait des méthodes économétriques adaptées (variables
instrumentales, discontinuités de régression).

\textbf{Qualité vs.~quantité.} Le niveau d'éducation atteint ne
renseigne pas sur la qualité des apprentissages, qui peut varier
fortement selon l'établissement, la région et la période de
scolarisation.

\begin{center}\rule{0.5\linewidth}{0.5pt}\end{center}

\section{Conclusion}\label{conclusion}

Cette analyse des données du Nigeria GHS Wave 4 met en lumière quatre
résultats structurants sur les inégalités éducatives au Nigeria.

Premièrement, la \textbf{distribution globale} révèle un déficit
significatif d'instruction formelle dans la population, avec une part
non négligeable d'individus n'ayant jamais fréquenté l'école.

Deuxièmement, les \textbf{inégalités de genre} sont statistiquement
significatives et économiquement substantielles, reflétant des normes
socioculturelles et des contraintes économiques qui ont historiquement
limité l'accès des femmes à l'éducation.

Troisièmement, le \textbf{gradient générationnel} constitue le résultat
le plus marquant : le test de Kruskal-Wallis (χ² = 287.28 ; p
\textless{} 0,001) et les comparaisons post-hoc de Dunn confirment que
chaque génération est significativement plus instruite que la
précédente, témoignant de l'impact réel des politiques d'éducation
universelle depuis les années 1970.

Quatrièmement, les \textbf{fractures spatiales} --- entre milieu rural
et urbain, et entre États du Nord et du Sud --- confirment la
persistance d'inégalités géographiques profondes qui se superposent aux
inégalités de genre et de génération.

Ces résultats plaident en faveur de politiques éducatives différenciées,
ciblant en priorité les femmes des zones rurales et les États du Nord,
et articulant des interventions sur l'offre (infrastructures,
enseignants) et la demande (transferts monétaires conditionnels,
sensibilisation).

\begin{center}\rule{0.5\linewidth}{0.5pt}\end{center}

\emph{Packages R utilisés : \texttt{haven}, \texttt{dplyr},
\texttt{ggplot2}, \texttt{scales}, \texttt{dunn.test}, \texttt{knitr},
\texttt{kableExtra}.}

\end{document}
